\chapter{Emissions and Environmental Markets}\label{ch:m3:emissions-and-environmental-markets}

By 30 September each year, a steel plant in Europe must hand over to the authorities one
allowance for every tonne of carbon dioxide it emitted the year before. On 21 February
2023 an allowance changed hands above 100 euros for the first time. For an industry
that emits tonnes of carbon dioxide for every tonne of product, allowances had become a
cost to hedge like its coal and ore. \omterm{def:m3:emissions-and-environmental-markets:cap}{Emission allowances} are a market created by law: the regulator fixes the quantity,
the market finds the price. This chapter describes how \omterm{def:m3:emissions-and-environmental-markets:cap}{cap-and-trade} works, how the
European system manages its supply, how carbon enters power prices through the \omterm{def:m3:power-markets-design:merit}{merit
order}, and the certificates that trade alongside.

\section{Cap and trade}

\begin{definition}[Cap-and-trade, emission allowance]\label{def:m3:emissions-and-environmental-markets:cap}
\emph{Cap-and-trade}\index{cap-and-trade} is a regulation that caps the total emissions of
covered installations over a period, issues that many \emph{emission
allowances}\index{emission allowance} (each a right to emit one tonne of CO$_2$ or its
equivalent), and requires every installation to hold and give up allowances for its
emissions, leaving the allowances free to be traded.
\end{definition}

The trade is the point: an installation that can cut a tonne for less than the allowance
price cuts it and sells the allowance; one that cannot buys. The cap is met by the cheapest
reductions available, whoever owns them, and the price tells everyone what a tonne is
worth. A tax fixes the price and lets the quantity follow; a cap fixes the quantity and lets
the price follow.

\begin{definition}[Allowance surrender, free allocation]\label{def:m3:emissions-and-environmental-markets:surrender}
\emph{Allowance surrender}\index{allowance surrender} is the annual handing-over by an
operator to its registry of allowances equal to its verified emissions of the previous year;
the surrendered allowances are cancelled. \emph{Free allocation}\index{free allocation} is the
distribution of part of the cap to installations without charge, by benchmark rules; the rest
is auctioned.
\end{definition}

\begin{omfigure}
\centering
\begin{tikzpicture}[font=\small]
\draw[thick,-{Stealth[length=2.5mm]}] (0,0) -- (13,0);
\foreach \x/\t/\l in {0.6/{1 Jan, year $t$}/{emissions\\begin},4.4/{31 Dec, year $t$}/{emissions of\\year $t$ end},7.6/{31 Mar, $t+1$}/{verified report\\submitted},11.4/{30 Sep, $t+1$}/{allowances\\surrendered}}{
  \draw[thick] (\x,0.12) -- (\x,-0.12);
  \node[anchor=south,align=center] at (\x,0.2) {\t};
  \node[anchor=north,align=center] at (\x,-0.2) {\l};}
\node[omDef,align=center,font=\footnotesize] at (2.5,-1.6) {free allocation and auctions\\throughout the year};
\node[omThm,align=center,font=\footnotesize] at (9.5,-1.6) {a shortfall is bought\\before the deadline};
\end{tikzpicture}

\omcaption{The compliance cycle of the European system for one year of emissions, with the
deadlines in force since 2024. Schematic.}\label{fig:m3:emissions-and-environmental-markets:cycle}
\end{omfigure}

\section{The European system}

The European Union's emissions trading system started in 2005 and covers power generation,
heavy industry, aviation within Europe and, since 2024, shipping. Its cap falls every year;
power generators buy all their allowances, while industries exposed to foreign competition
receive part of theirs free. Allowances not allocated free are sold at regular auctions, and they trade as spot and as
futures.

\begin{dated}{2026-09}{European compliance and scope}\label{dat:m3:emissions-and-environmental-markets:eu}
Since 2024 operators surrender allowances for the previous year's emissions by 30 September
(before, 30 April), after submitting a verified emissions report by 31 March. An operator short on the day
pays an excess emissions penalty of EUR~100 a tonne, indexed to consumer prices for allowances issued from 2013,
and must still surrender the missing allowances the following year. A second system
for fuels used in buildings and road transport (ETS2) was due to start in 2027; in March 2026
the Parliament and the Council agreed to postpone it to 2028.
\end{dated}

\section{The market stability reserve}

A cap fixes the quantity, but when demand falls (a recession, a pandemic, cheap renewables) a
fixed quantity leaves a surplus that depresses the price for years. The European system
answers with a rule that adjusts future auction volumes.

\begin{definition}[Market stability reserve, total number of allowances in circulation]\label{def:m3:emissions-and-environmental-markets:msr}
The \emph{total number of allowances in circulation}\index{total number of allowances in circulation}
(TNAC) is the surplus of allowances held in the market: all allowances issued and not
cancelled, minus those surrendered for emissions and those held in the reserve, published each
year by the Commission. The \emph{market stability reserve}\index{market stability reserve}
(MSR) is a reserve into which allowances are withheld from future auctions when the TNAC is
high, and from which they are released when it is low.
\end{definition}

\begin{proposition}[The intake rule]\label{prop:m3:emissions-and-environmental-markets:msr}
Under the rules in force since 2024, the reserve takes over the twelve months from September:
24\% of the TNAC if the TNAC exceeds 1\,096 million; the TNAC minus 833 million if it lies
between 833 and 1\,096 million; nothing between 400 and 833 million; and it releases 100
million if the TNAC falls below 400 million. Allowances held in the reserve above 400 million
cease to be valid.
\end{proposition}

\begin{proof}
This is the rule as amended in 2023; the continuous piece between 833 and 1\,096 million
removes the jump that a flat 24\% would cause at 833 million
(\cref{fig:m3:emissions-and-environmental-markets:msr}).
\end{proof}

\begin{example}[Two published years]\label{ex:m3:emissions-and-environmental-markets:msr}
The TNAC published in 2025 was 1\,148\,049\,585 allowances, above 1\,096 million: the
reserve takes 24\%, 275\,531\,900 allowances, from September 2025 to August 2026. The TNAC
published on 29 May 2026 was 1\,023\,494\,202: the reserve takes the excess over 833 million,
190\,494\,202 allowances, from September 2026 to August 2027.
\end{example}

\begin{omfigure}
\begin{tikzpicture}
\begin{axis}[width=11cm,height=5cm,
  xlabel={TNAC, million allowances},ylabel={intake, million},
  tick label style={font=\small},label style={font=\small},
  xmin=200,xmax=1500,ymin=-150,ymax=400,grid=major,grid style={gray!25},
  xtick={400,833,1096,1400},
  legend columns=2,legend style={font=\footnotesize,at={(0.5,-0.3)},anchor=north,draw=none,fill=none,/tikz/every even column/.append style={column sep=8pt}},
  table/col sep=comma]
\addplot[omDef,very thick] table[x=tnac,y=intake]{figdata/markets-3/07-emissions-and-environmental-markets/msr.csv};
\addplot[omThm,only marks,mark=*,mark size=2.2pt] table[x=tnac,y=intake]{figdata/markets-3/07-emissions-and-environmental-markets/msrpoints.csv};
\draw[gray] (axis cs:200,0) -- (axis cs:1500,0);
\legend{intake rule,TNAC published in 2025 and 2026}
\end{axis}
\end{tikzpicture}

\omcaption{Allowances withheld from auctions over twelve months as a function of the published
surplus (negative: released), with the two most recent publications. Data: the Decision as
amended in 2023; Commission publications of 2025 and 2026.}\label{fig:m3:emissions-and-environmental-markets:msr}
\end{omfigure}

The reserve turns a fixed cap into a supply curve: the more the market holds, the fewer
allowances are auctioned. A trader who forecasts the TNAC forecasts next year's auction supply.

\section{Carbon in the merit order}

A plant's marginal cost includes its allowances: $(P_{\mathrm{fuel}} + e\,P_{\mathrm{CO_2}})/\eta$
(\cref{ch:m3:power-markets-design}). Coal emits more per MWh of fuel than gas and burns it less
efficiently, so a higher carbon price moves coal plants up the \omterm{def:m3:power-markets-design:merit}{merit order}.

\begin{definition}[Clean spark spread, clean dark spread]\label{def:m3:emissions-and-environmental-markets:clean}
The \emph{clean spark spread}\index{clean spark spread} is the \omterm{def:m3:power-markets-design:spark}{spark spread} minus the cost of the
allowances for one MWh of electricity from a gas plant, $P_{\mathrm{power}} - (P_{\mathrm{gas}} +
e_g P_{\mathrm{CO_2}})/\eta_g$; the \emph{clean dark spread}\index{clean dark spread} is the same
for a coal plant.
\end{definition}

\begin{definition}[Switching price]\label{def:m3:emissions-and-environmental-markets:switch}
The coal-to-gas \emph{switching price}\index{switching price} is the allowance price at which a
given gas plant and a given coal plant have the same marginal cost:
\[
P^{\star}_{\mathrm{CO_2}} = \frac{P_{\mathrm{gas}}/\eta_g - P_{\mathrm{coal}}/\eta_c}{e_c/\eta_c - e_g/\eta_g}.
\]
\end{definition}

\begin{example}[Where gas beats coal]\label{ex:m3:emissions-and-environmental-markets:switch}
With gas at \qty{35}{EUR/MWh} of fuel, coal at \qty{12}{}, efficiencies of 55\% and 40\%, and
emission factors of 0.202 and 0.341 tonnes per MWh of fuel (illustrative round numbers), a MWh of
electricity emits 0.367~t from gas and 0.853~t from coal, and the \omterm{def:m3:emissions-and-environmental-markets:switch}{switching price} is
\qty{69.32}{EUR/t}. At \qty{70}{EUR/t} and power at \qty{100}{EUR/MWh} the \omterm{def:m3:emissions-and-environmental-markets:clean}{clean spark spread} is
\qty{10.65}{} and the \omterm{def:m3:emissions-and-environmental-markets:clean}{clean dark spread} \qty{10.32}{EUR/MWh}: gas is just ahead.
\end{example}

\begin{omfigure}
\begin{tikzpicture}
\begin{axis}[width=11cm,height=4.8cm,
  xlabel={allowance price, EUR/t},ylabel={EUR/MWh},
  tick label style={font=\small},label style={font=\small},
  xmin=0,xmax=150,ymin=-60,ymax=80,grid=major,grid style={gray!25},
  legend columns=2,legend style={font=\footnotesize,at={(0.5,-0.3)},anchor=north,draw=none,fill=none,/tikz/every even column/.append style={column sep=8pt}},
  table/col sep=comma]
\addplot[omDef,very thick] table[x=carbon,y=spark]{figdata/markets-3/07-emissions-and-environmental-markets/spreads.csv};
\addplot[omThm,very thick,dashed] table[x=carbon,y=dark]{figdata/markets-3/07-emissions-and-environmental-markets/spreads.csv};
\draw[gray] (axis cs:69.32,-60) -- (axis cs:69.32,80);
\node[font=\footnotesize,anchor=south west] at (axis cs:70.5,40) {switching price};
\legend{clean spark spread,clean dark spread}
\end{axis}
\end{tikzpicture}

\omcaption{Clean spreads of the plants of \cref{ex:m3:emissions-and-environmental-markets:switch}
against the allowance price, with power at \qty{100}{EUR/MWh}. Coal's spread falls 2.3 times as
fast; the lines cross at the \omterm{def:m3:emissions-and-environmental-markets:switch}{switching price}. Data: the chapter's
tutorial.}\label{fig:m3:emissions-and-environmental-markets:spreads}
\end{omfigure}

The \omterm{def:m3:emissions-and-environmental-markets:switch}{switching price} moves with gas. With gas at \qty{20}{EUR/MWh} of fuel it is
\qty{13.11}{EUR/t}; with gas at \qty{50}{EUR/MWh} it is \qty{125.53}{EUR/t}
(\cref{fig:m3:emissions-and-environmental-markets:switch}). With European gas between \$27 and
\$70 per MMBtu in every month of 2022 (\cref{ch:m3:natural-gas-and-lng}), the \omterm{def:m3:emissions-and-environmental-markets:switch}{switching price} was
far above any allowance price seen, and coal ran ahead of gas.

\begin{omfigure}
\begin{tikzpicture}
\begin{axis}[width=11cm,height=4.4cm,
  xlabel={gas price, EUR/MWh of fuel},ylabel={switching price, EUR/t},
  tick label style={font=\small},label style={font=\small},
  xmin=10,xmax=60,ymin=0,ymax=180,ytick={50,100,150},grid=major,grid style={gray!25},
  table/col sep=comma]
\addplot[omProp,very thick] table[x=gas,y=carbon]{figdata/markets-3/07-emissions-and-environmental-markets/switch.csv};
\node[font=\footnotesize,anchor=north west,align=left] at (axis cs:12,170) {above the line: gas cheaper\\below: coal cheaper};
\end{axis}
\end{tikzpicture}

\omcaption{The allowance price that makes the gas and coal plants of
\cref{ex:m3:emissions-and-environmental-markets:switch} equal, as the gas price varies with coal at
\qty{12}{EUR/MWh}. Data: the chapter's tutorial.}\label{fig:m3:emissions-and-environmental-markets:switch}
\end{omfigure}

\section{Other systems and the border adjustment}

\begin{dated}{2026-09}{Emissions trading worldwide}\label{dat:m3:emissions-and-environmental-markets:world}
The International Carbon Action Partnership counted 41 emissions trading systems in force in
February 2026, covering 26\% of global greenhouse-gas emissions in jurisdictions that produce 63\%
of world GDP, with national systems launching in 2026 in Japan, India and Vietnam.
\end{dated}

A system that prices carbon at home but not on imports invites production to move abroad. The
European answer is a charge at the border.

\begin{definition}[Carbon border adjustment mechanism]\label{def:m3:emissions-and-environmental-markets:cbam}
The \emph{carbon border adjustment mechanism}\index{carbon border adjustment mechanism} (CBAM)
requires importers of certain carbon-intensive goods into the European Union (iron and steel,
aluminium, cement, fertilisers, electricity, hydrogen) to buy certificates for the emissions embedded
in them, priced on the allowance price, less any carbon price paid where they were made.
\end{definition}

\begin{dated}{2026-09}{CBAM timetable}\label{dat:m3:emissions-and-environmental-markets:cbam}
The definitive period began on 1 January 2026. Regulation (EU) 2025/2083 exempts importers of less
than 50 tonnes a year of covered goods (other than hydrogen and electricity) and postponed the start
of certificate sales from 1 January 2026 to 1 February 2027.
\end{dated}

\section{Certificates of origin and carbon credits}

Allowances are one of several environmental instruments. Others certify attributes rather than
rights to emit.

\begin{definition}[Guarantee of origin, renewable energy certificate]\label{def:m3:emissions-and-environmental-markets:go}
A \emph{guarantee of origin}\index{guarantee of origin} is an electronic document, issued in the
European Union for each MWh of renewable energy on request of its producer, whose sole function is
to prove to a final customer that a given quantity of energy was produced from renewable sources.
A \emph{renewable energy certificate}\index{renewable energy certificate} is the US equivalent: a
tradable instrument representing the renewable attributes of one MWh of renewable electricity.
\end{definition}

\begin{definition}[Carbon credit]\label{def:m3:emissions-and-environmental-markets:credit}
A \emph{carbon credit}\index{carbon credit} (or offset) is a certificate representing one tonne of
CO$_2$ equivalent reduced or removed by a project outside a cap, issued by a crediting standard, and
used voluntarily by buyers to claim compensation for their own emissions.
\end{definition}

The three instruments differ in what backs them. An allowance is a legal right within a cap, its
supply fixed by law. A \omterm{def:m3:emissions-and-environmental-markets:go}{guarantee of origin} certifies how a MWh was produced; its price reflects the
demand of consumers who want green claims. A \omterm{def:m3:emissions-and-environmental-markets:credit}{carbon credit} is a claim about a counterfactual (what
would have been emitted without the project), and its value rests on the credibility of that claim.

\section{Tutorial: the reserve and the switching price}\label{tut:m3:emissions-and-environmental-markets:switch}

\begin{tutorial}
\textbf{Goal.} Apply the reserve's rule to the published surplus, and compute the \omterm{def:m3:emissions-and-environmental-markets:switch}{switching price}
and clean spreads. \textbf{End state:}
\cref{ex:m3:emissions-and-environmental-markets:msr,ex:m3:emissions-and-environmental-markets:switch}
and \cref{fig:m3:emissions-and-environmental-markets:msr,fig:m3:emissions-and-environmental-markets:spreads}.
\begin{tutsteps}
\item \textbf{The rule.} \Cref{prop:m3:emissions-and-environmental-markets:msr}.
\omcode{code/firm/carbon/firm_carbon.py}{12}{23}{The reserve's intake rule.}\label{lst:m3:emissions-and-environmental-markets:msr}
\item \textbf{Spreads and switching.}
\omcode{code/firm/carbon/firm_carbon.py}{26}{40}{Emissions per MWh, clean spreads and the switching price.}\label{lst:m3:emissions-and-environmental-markets:switch}
\item \textbf{Run} \texttt{m3\_carbon.msr\_cases()}, \texttt{utility()} and \texttt{fig\_carbon.py}.
\end{tutsteps}
\textbf{What to change next.} Lower coal's efficiency to 36\% and find the new \omterm{def:m3:emissions-and-environmental-markets:switch}{switching price}; model
a year in which the TNAC falls below 833 million and the reserve stops taking allowances.
\end{tutorial}

\section{Build: the compliance book}\label{bld:m3:emissions-and-environmental-markets:carbon}

\begin{build}
\bfield{Purpose} The miniature firm's generating assets need allowances; its traders price clean
spreads and forecast auction supply. The compliance book keeps allowances against emissions and
computes what the reserve will do next September.
\bfield{Interface} \texttt{msr\_intake(tnac)}; \texttt{emissions\_per\_mwh(efficiency, factor)};
\texttt{clean\_spread(power, fuel, carbon, efficiency, factor)}; \texttt{switching\_price(gas, coal,
eff\_gas, eff\_coal)}; \texttt{ComplianceAccount.buy / emit / surrender}.
\bfield{Rules} Tonnes are integers; surrender covers the oldest uncovered emissions and returns the
shortfall; the reserve rule reproduces the Commission's published intakes to the allowance.
\bfield{Acceptance tests} \texttt{code/firm/carbon/tests/}: the 2025 and 2026 intakes exactly; equal
clean spreads at the \omterm{def:m3:emissions-and-environmental-markets:switch}{switching price}; a shortfall and its cure.
\bfield{Stretch} Read the day-ahead prices of \texttt{firm.dayahead} and compute each plant's clean
spread per hour; a forward curve of allowances with the cost of carry.
\end{build}

\begin{omsources}
\item European Commission, Market Stability Reserve page and the Communications on the TNAC of 2025
and 2026; Decision (EU) 2015/1814 as amended by Decision (EU) 2023/852.
\item Directive (EU) 2023/959 (surrender deadline); Directive 2003/87/EC, consolidated text of 1 March 2024, Article 16(3)--(4)
(Internet Archive copy of EUR-Lex); Council and Parliament agreement on ETS2 (2026).
\item Regulation (EU) 2023/956 (CBAM) and Regulation (EU) 2025/2083.
\item Directive (EU) 2018/2001, Article 19 (guarantees of origin).
\item ICAP, \emph{Emissions Trading Worldwide: Status Report 2026}.
\item Reuters report of 21 February 2023 on the EUA price (via MarketScreener).
\end{omsources}

\section{Exercises}

\begin{exercise}[$\star$]\label{exo:m3:emissions-and-environmental-markets:1}
A plant emitted 1\,000 tonnes, holds 900 allowances and buys 100 more before the deadline. What
does it surrender, and when?
\end{exercise}

\begin{exercise}[$\star$]\label{exo:m3:emissions-and-environmental-markets:2}
What does the reserve do if the published TNAC is 1\,200 million? 900 million? 350 million?
\end{exercise}

\begin{exercise}[$\star$]\label{exo:m3:emissions-and-environmental-markets:3}
How does a cap differ from a tax when demand for emitting falls sharply?
\end{exercise}

\begin{exercise}[$\star\star$]\label{exo:m3:emissions-and-environmental-markets:4}
With the plants of \cref{ex:m3:emissions-and-environmental-markets:switch}, what is the \omterm{def:m3:emissions-and-environmental-markets:switch}{switching
price} if gas costs \qty{20}{EUR/MWh}? And if coal's efficiency is 38\%, with gas at \qty{35}{}?
\end{exercise}

\begin{exercise}[$\star\star$]\label{exo:m3:emissions-and-environmental-markets:5}
How many tonnes does a 55\%-efficient gas plant emit per MWh of electricity, and what does a 10-euro
rise in the allowance price add to its marginal cost?
\end{exercise}

\begin{exercise}[$\star\star$]\label{exo:m3:emissions-and-environmental-markets:6}
Why does a \omterm{def:m3:emissions-and-environmental-markets:go}{guarantee of origin} not reduce anyone's emissions directly, while a surrendered allowance
does?
\end{exercise}

\begin{exercise}[$\star\star\star$]\label{exo:m3:emissions-and-environmental-markets:7}
\emph{Coding.} With \texttt{msr\_intake}, find the TNAC between 800 and 1\,200 million at which the
intake is largest relative to the TNAC, and show that the rule is continuous at 1\,096 million
within a million.
\end{exercise}

\begin{exercise}[$\star\star\star$]\label{exo:m3:emissions-and-environmental-markets:8}
\emph{Find the flaw.} ``Allowances are cheap this year, so the cap is not binding and has no effect
on emissions.''
\end{exercise}

\section{Problem: The Utility's Carbon Hedge}

\begin{problem}[{Weekend problem --- coal, gas and next year's allowances}]\label{pb:m3:emissions-and-environmental-markets:1}
A utility plans to generate 6~TWh from coal (40\% efficiency) and 4~TWh from gas (55\%) next year,
with the fuels and emission factors of \cref{ex:m3:emissions-and-environmental-markets:switch}. It
receives no \omterm{def:m3:emissions-and-environmental-markets:surrender}{free allocation}. Allowances trade at \qty{70}{EUR/t}.

\medskip\noindent\textbf{Part I --- Needs.}
\begin{enumerate}
\item How many tonnes does each MWh from coal and from gas emit?
\item How many allowances does the plan need?
\item What do they cost at \qty{70}{EUR/t}?
\item When must they be surrendered?
\item Why does the utility buy them forward rather than when it emits?
\end{enumerate}

\medskip\noindent\textbf{Part II --- Switching.}
\begin{enumerate}[resume]
\item Give the \omterm{def:m3:emissions-and-environmental-markets:switch}{switching price}.
\item At \qty{70}{EUR/t}, which plant is cheaper, and by how much per MWh?
\item How many tonnes would moving 1~TWh from coal to gas save?
\item At what allowance price does that move pay?
\item What happens to the \omterm{def:m3:emissions-and-environmental-markets:switch}{switching price} if gas rises to \qty{50}{EUR/MWh}?
\end{enumerate}

\medskip\noindent\textbf{Part III --- The market.}
\begin{enumerate}[resume]
\item What does the reserve's rule imply for next year's auction supply, given the TNAC of 2026?
\item How would a fall in industrial output affect the allowance price and the TNAC?
\item Why do power prices rise with allowance prices in a coal-marginal hour?
\item Who holds allowances besides compliance buyers, and why?
\item What does the border adjustment change for the utility?
\end{enumerate}

\medskip\noindent\textbf{Part IV --- Judgement.}
\begin{enumerate}[resume]
\item Why might the utility hedge its power sales and its allowances together?
\item What does a \omterm{def:m3:emissions-and-environmental-markets:clean}{clean dark spread} near zero say about the coal plant's future?
\item Is the \omterm{def:m3:emissions-and-environmental-markets:switch}{switching price} a floor or a ceiling for allowances in a gas-rich year?
\item State the \emph{named result}: the \omterm{def:m3:emissions-and-environmental-markets:switch}{switching price} and the allowances the plan needs.
\item In one sentence: what sets the price of an allowance?
\end{enumerate}
\end{problem}

\section{Interview questions}

\begin{interviewq}[$\star$ \iqroles{trader}]\label{iq:m3:emissions-and-environmental-markets:1}
What is the difference between a carbon tax and \omterm{def:m3:emissions-and-environmental-markets:cap}{cap-and-trade}?
\end{interviewq}

\begin{interviewq}[$\star$ \iqroles{trader, researcher}]\label{iq:m3:emissions-and-environmental-markets:2}
Define the \omterm{def:m3:emissions-and-environmental-markets:clean}{clean spark spread}. When is it negative?
\end{interviewq}

\begin{interviewq}[$\star\star$ \iqroles{researcher}]\label{iq:m3:emissions-and-environmental-markets:3}
How does the \omterm{def:m3:emissions-and-environmental-markets:msr}{market stability reserve} change the elasticity of allowance supply?
\end{interviewq}

\begin{interviewq}[$\star\star$ \iqroles{trader}]\label{iq:m3:emissions-and-environmental-markets:4}
Gas prices double. What happens to allowance prices, and why?
\end{interviewq}

\begin{interviewq}[$\star\star$ \iqroles{risk}]\label{iq:m3:emissions-and-environmental-markets:5}
A company buys \omterm{def:m3:emissions-and-environmental-markets:credit}{carbon credits} to claim neutrality. What risks does it carry?
\end{interviewq}

\begin{interviewq}[$\star\star\star$ \iqroles{researcher, developer}]\label{iq:m3:emissions-and-environmental-markets:6}
Build a model of next year's allowance demand from the power sector. What data and what structure?
\end{interviewq}
